\chapter{Introduction}
\label{chap:introduction}

\section{Background}
\label{section:background}

Meetings are an indispensable tool for collaborative work in both academic and professional settings.
In theory, a focused meeting covering a single topic should require no more than 15--30 minutes.
In practice, however, teams routinely turn short discussions into multi-hour sessions.
The root cause is \textit{topic drift}: the gradual, often unnoticed, departure of conversation from the
original agenda.
 
The development team behind MeetRX experienced this problem first-hand during group project work.
What should have been a concise weekly sync repeatedly stretched into a three-hour session because
members lost track of the intended focus and the conversation wandered across unrelated topics.
This observation motivated the creation of an intelligent meeting assistant that can detect drift as it
occurs and help teams stay aligned.

\section{Problem Statement}
\label{section:problem-statement}

The following problems were identified through the team's own experience and stakeholder interviews:
 
\begin{enumerate}
    \item \textbf{Burden on the minute-taker.} During meetings, the person responsible for taking notes
    cannot participate meaningfully in the discussion. Their attention is divided between listening,
    comprehending, and writing, which reduces both the quality of the notes and their contribution to
    the meeting.
 
    \item \textbf{Standup meetings overrunning their allotted time.} Daily standups and progress-update
    meetings are intended to be short, time-boxed events. In practice, they frequently exceed their
    scheduled duration because participants introduce tangential topics, triggering lengthy discussions
    that should be handled separately.
 
    \item \textbf{Repeated debates caused by poor institutional memory.} In long-term projects with
    frequent meetings, teams often forget what was discussed or decided in earlier sessions. As a result,
    the same topics are re-opened and the same decisions are re-debated, consuming time that could be
    spent on new work.
\end{enumerate}

\section{Solution Overview}
\label{section:solution-overview}

MeetRX addresses the problems above through three integrated AI subsystems:
 
\begin{enumerate}
    \item \textbf{Summarisation AI}: Automatically transcribes meeting audio using OpenAI Whisper
    and generates a structured summary so that no human minute-taker is required. All participants can
    engage fully in the discussion.
 
    \item \textbf{Topic Drift Detection AI (Main Feature)}: Continuously analyses the live conversation
    using natural language processing techniques to assign a drift score to recent utterances. When the
    score exceeds a configurable threshold, a large language model re-evaluation is triggered and a
    real-time alert is surfaced to meeting participants, prompting them to return to the agenda.
 
    \item \textbf{Agentic AI}: Provides a conversational interface that can answer questions about past
    meetings, look up decisions recorded in previous summaries, query the project codebase, and integrate
    with external tools to retrieve relevant context on demand.
\end{enumerate}

\subsection{Features}
\label{subsection:features}

\begin{enumerate}[leftmargin=80pt]
    \item \textbf{Real-time Topic Drift Detection}: Monitors conversation flow and raises an alert when
    discussion deviates from the stated agenda topic.
    \item \textbf{Automatic Meeting Summarisation}: Produces a concise, structured summary at the end of
    every session without requiring manual note-taking.
    \item \textbf{Q\&A Chatbot}: Allows participants to query past meeting summaries and project context
    using natural language.
    \item \textbf{Task Assignment}: Extracts and records action items from meeting transcripts and assigns
    them to the appropriate team members.
\end{enumerate}

\section{Target User}
\label{section:target-user}

MeetRX is designed for three primary user groups:
 
\begin{itemize}
    \item \textbf{Students}: Particularly those conducting group project meetings who need help keeping
    discussions on-topic and producing reliable records of decisions made.
    \item \textbf{Software Developers}: Engineering teams running daily standups, sprint planning, and
    retrospectives who want shorter, more focused meetings and automatic summaries.
    \item \textbf{Scrum Masters}: Facilitators responsible for sprint ceremonies who need tooling to
    enforce time-boxes and surface relevant historical context during discussions.
\end{itemize}
 
\section{Terminology}
\label{section:terminology}

\begin{description}
    \item[Topic Drift] A situation in which the ongoing conversation is no longer contributing to the
    intended meeting topic or agenda item.
    \item[Drift Score] A numerical value computed by the NLP pipeline that quantifies how far the current
    utterance has deviated from the baseline topic embedding.
    \item[Agenda] A predefined list of topics that a meeting is intended to cover.
    \item[Utterance] A single spoken turn or message produced by one participant during a meeting.
    \item[Summarisation] The process of condensing a meeting transcript into a shorter document that
    captures the key points, decisions, and action items.
    \item[Agentic AI] An AI system capable of autonomously invoking external tools or data sources to
    answer queries without explicit step-by-step instructions from the user.
    \item[LLM] Large Language Model; a neural network trained on large text corpora capable of generating
    and understanding natural language.
    \item[Whisper] An open-source automatic speech recognition model developed by OpenAI.
\end{description}